\chapter{תקציר מנהלים}\label{ch:executive-summary}

\noindent
\keynumber{\RRrevFinal{}M \shekel}{הכנסות ב-2031}\hfill
\keynumber{\RRroofsFinal}{גגות פעילים ב-2031}\hfill
\keynumber{\RRroofPayback{} שנים}{החזר השקעה לגג}

\bigskip

\section*{הבעיה}
בערים הצפופות של גוש דן גגות רבים של בנייני מגורים משמשים לדודי שמש ולציוד בלבד, ושום דבר לא גדל עליהם. במקביל, ירקות העלים של המסעדות מגיעים מרחוק, אחרי ימים בדרך, במדינה שבה מים הם משאב מוגבל. הגג הוא נכס ריק, והירק הוא מוצר מתכלה שמאבד טריות בהובלה.

\section*{הפתרון}
גגון (\en{RoofRoot}) מתקינה מערכות הידרופוניקה מסוג \en{NFT} על גגות של בנייני מגורים בני 6 קומות ומעלה, במסגרת הסכם ל-\RRcontractYears{} שנים עם ועד הבית. הוועד מקבל \RRrentShare\% מהכנסות הגג. את המערכת מנהל \en{RootOS}: סוכן בינה מלאכותית שקורא חיישנים, מפעיל משאבות והצללה, ומעביר פעולות חריגות לאישור מגדל אנושי.

\section*{השוק והמודל העסקי}
הלקוחות הם מסעדות ובתי קפה ברדיוס של ק״מ אחד מהגג (\en{B2B}), מנוי ״סל ירוק״ לדיירי הבניין, ומ-2029 חוות חיצוניות שיעשו שימוש ברישיון \en{RootOS}. לפי הערכת צוות, שוק יעד מצרפי של \RRtam{} מיליון \shekel, שוק ניתן לשירות של \RRsam{} מיליון \shekel, ויעד נתח של \RRsom{} מיליון \shekel{} ב-2031 (פירוט בפרק~\ref{ch:market}).

\section*{כלכלת הגג}
כל גג כולל \RRarea{} מ״ר שטח גידול ומניב \RRyield{} ק״ג/מ״ר בשנה, כלומר \RRroofKg{} ק״ג בשנה. במחיר ממוצע של \RRprice{} \shekel{} לק״ג, ההכנסה היא \RRroofRevenue{} אלף \shekel{} בשנה. ההשקעה בגג היא \RRcapexRoof{} אלף \shekel, והחזרה תוך \RRroofPayback{} שנים.

\section*{התוכנית וההישגים הצפויים}
אנו מתכננים \RRroofsFirst{} גגות ב-2027 ו-\RRroofsFinal{} גגות ב-2031. התוכנית היא תוכנית ולא ביצוע: אין עדיין הכנסות בפועל, ולכן הפרק על הסיכונים (פרק~\ref{ch:risks}) מציג מה צריך להצליח כדי שהמספרים יתממשו.

\section*{למה גגון תצליח}
שלושה יתרונות: נכס זול (גג שהוועד מעוניין לנצל), ביקוש מקומי צפוי של מסעדות סמוכות (טרם אומת, ייבדק בפיילוט), ושכבת תוכנה (\en{RootOS}) שמקטינה את עלות העבודה לגג ומאפשרת למגדל אחד לנהל גגות רבים. הצוות המייסד משלב ידע אגרונומי (נועה ברק, מנכ״לית), למידת מכונה ורובוטיקה (עידו שחר, מנכ״ל טכנולוגיות) ותפעול לוגיסטי (מאיה לוי, מנהלת תפעול).

\section*{הכותרת הפיננסית}
הכנסות של \RRrevFinal{} מיליון \shekel{} ב-2031, רווח \en{EBITDA} של \RRebitdaFinal{} מיליון \shekel{} (שיעור של \RRebitdaMarginFinal\%) ורווח גולמי של \RRgrossMarginFinal\%. החברה הופכת ל-\en{EBITDA} חיובי ב-\RRbreakEvenYear. הערך הנוכחי נקי (\en{NPV}) הוא \RRnpv{} מיליון \shekel{} בשיעור היוון של \RRdiscount\%. יתרת המזומנים הנמוכה ביותר היא \RRminCash{} מיליון \shekel{} ב-\RRminCashYear{} (ראו איור~\ref{fig:cash-curve}).

\section*{הבקשה}
אנו מגייסים \RRseed{} מיליון \shekel{} בסבב \en{Seed} ב-2027, וסבב \en{Series A} של \RRseriesA{} מיליון \shekel{} מתוכנן ל-2029; סך ההון הנדרש הוא \RRequityTotal{} מיליון \shekel. פירוט השימוש בכספים נמצא בטבלה~\ref{tab:use-of-funds}.

\begin{insightbox}
גג אחד מחזיר את ההשקעה בו תוך \RRroofPayback{} שנים. הסיכון המרכזי אינו הכלכלה של גג בודד אלא קצב ההתקנה והיכולת להפעיל עשרות גגות עם צוות קטן; \en{RootOS} נועד לפתור בדיוק את זה.
\end{insightbox}
